\documentclass[11pt]{article}

\usepackage[utf8]{inputenc}
\usepackage[T1]{fontenc}
\usepackage{geometry}
\usepackage{amsmath}
\usepackage{booktabs}
\usepackage{array}

\geometry{letterpaper, margin=2.5cm}

\title{Informe del módulo de interfaz visual}
\author{CSDI RAG Engine}
\date{}

\begin{document}

\maketitle

\section{Módulo de interfaz visual}

El módulo de interfaz visual es la capa de presentación del sistema CSDI RAG Engine. Su
responsabilidad es exponer todas las capacidades del backend a través de una aplicación web
interactiva que un operador o usuario final puede utilizar desde un navegador, sin necesidad
de interactuar directamente con la API REST ni con herramientas de línea de comandos.

La interfaz no implementa lógica de negocio: actúa como cliente de la API del servidor, que
opera en el puerto configurado por la variable \texttt{APP\_PORT} (por defecto \texttt{8888}).
Todos los datos que muestra provienen de respuestas de esa API; no existe estado local
persistente en el cliente más allá del identificador de sesión de conversación. Esta separación
es deliberada: mantener la interfaz como cliente puro de la API garantiza que cualquier
modificación al pipeline, al índice o a los modelos sea inmediatamente reflejada en lo que
muestra la interfaz, sin necesidad de coordinar cambios entre capas.

La interfaz está compuesta por ocho secciones principales, cada una dedicada a un dominio
funcional distinto del sistema: consulta en lenguaje natural, historial de conversación,
búsqueda directa sobre el índice, configuración del pipeline, gestión de fuentes e ingesta,
monitorización operativa, retroalimentación de consultas y evaluación de calidad.

\section{Panel de chat RAG}

\paragraph{Función del panel.}
Es la sección de uso cotidiano del sistema. Permite formular preguntas en texto libre y recibir
respuestas generadas por el modelo de lenguaje, fundamentadas en fragmentos recuperados del
corpus de documentación interno o de la web. Es el punto de entrada principal para cualquier
usuario que desee consultar información del corpus sin conocer la estructura del índice.

\paragraph{Datos presentados.}
La respuesta del endpoint \texttt{POST /api/v1/rag/query} incluye todos los campos necesarios
para que la interfaz muestre tanto la respuesta como su procedencia completa. El campo
\textbf{answer} contiene el texto generado por el modelo en español. El campo \textbf{sources}
es una lista de fragmentos utilizados para construir la respuesta: cada fuente incluye el
título del documento, la URL de origen, el tipo de fuente (\texttt{corpus} o
\texttt{web\_cache}), el método de recuperación (\texttt{bm25}, \texttt{vector},
\texttt{hybrid} o \texttt{web\_cache}), la puntuación de relevancia, la puntuación de frescura,
la prioridad de visualización y el rango de posicionamiento. El campo \textbf{model} identifica
el modelo exacto que generó la respuesta. Los campos \textbf{prompt\_tokens} y
\textbf{completion\_tokens} reportan el consumo de tokens de la llamada. Cuando se ejecutó
búsqueda web, los campos \textbf{cache\_searched}, \textbf{cache\_hits},
\textbf{external\_search\_executed} y \textbf{external\_indexed\_count} informan sobre el
alcance de esa búsqueda.

La respuesta presenta una separación visual entre la información proveniente del corpus interno
y la obtenida de búsqueda web. Cuando todos los fragmentos son del corpus, la respuesta no
incluye ninguna indicación de búsqueda web. Cuando los fragmentos son exclusivamente de la web,
la respuesta indica explícitamente que no se encontró información en el corpus interno. Cuando
hay fragmentos de ambas fuentes, la respuesta distingue qué parte de la información proviene
de cada origen.

\paragraph{Parámetros editables.}

\begin{center}
\begin{tabular}{lp{9cm}}
\toprule
\textbf{Parámetro} & \textbf{Descripción} \\
\midrule
\texttt{query}      & Pregunta en texto libre. Longitud mínima: 1 carácter. \\
\texttt{session\_id} & Identificador de sesión opcional (8--128 caracteres alfanuméricos
                       y guiones). Cuando se incluye, el intercambio queda registrado en
                       el historial de conversación. \\
\bottomrule
\end{tabular}
\end{center}

\section{Panel de historial de conversación}

\paragraph{Función del panel.}
Muestra la secuencia completa de preguntas y respuestas de una sesión identificada por su
\texttt{session\_id}. Permite al usuario revisar interacciones anteriores sin necesidad de
repetir consultas. El historial es consultable en cualquier momento durante el tiempo de vida
de la sesión, que por defecto es de 7 días (\texttt{CHAT\_HISTORY\_TTL\_SECONDS=604800}),
renovado en cada intercambio.

\paragraph{Datos presentados.}
El endpoint \texttt{GET /api/v1/rag/history/\{session\_id\}} devuelve una lista ordenada de
mensajes. Cada entrada incluye un identificador único, el tipo de mensaje (\texttt{"user"} para
preguntas, \texttt{"assistant"} para respuestas), el contenido textual, la marca temporal en
formato ISO 8601, la lista de fuentes citadas y el nombre del modelo utilizado en el caso de
los mensajes del asistente.

\paragraph{Parámetros editables.}

\begin{center}
\begin{tabular}{lp{9cm}}
\toprule
\textbf{Parámetro} & \textbf{Descripción} \\
\midrule
\texttt{session\_id} & Identificador de sesión a consultar. Debe coincidir con el usado
                       durante las consultas en el panel de chat. \\
\bottomrule
\end{tabular}
\end{center}

\section{Panel de búsqueda híbrida}

\paragraph{Función del panel.}
Permite ejecutar búsquedas directas sobre el índice sin pasar por el modelo generativo.
Su utilidad principal es exploratoria y de diagnóstico: permite verificar qué fragmentos
recupera el sistema ante una consulta concreta, comparar los resultados de la búsqueda
léxica frente a la vectorial, y depurar el comportamiento del pipeline de recuperación.
Expone tres modalidades: búsqueda híbrida sobre el corpus, búsqueda híbrida sobre el caché
web y búsqueda con expansión por retroalimentación de consultas.

\paragraph{Datos presentados.}
El endpoint \texttt{POST /api/v1/search} devuelve por cada resultado el identificador del
fragmento (\texttt{chunk\_id}), la puntuación de relevancia RRF (\texttt{score}), la URL
del documento de origen, el título, el \emph{breadcrumb} de navegación, el texto completo
del fragmento y el conjunto \texttt{found\_by}, que indica si el fragmento fue encontrado
por BM25, por el recuperador vectorial o por ambos. La búsqueda sobre el caché web
(\texttt{POST /api/v1/search/web-cache}) devuelve la misma estructura pero limitada a los
fragmentos indexados de búsquedas web anteriores.

\paragraph{Parámetros editables.}

\begin{center}
\begin{tabular}{lp{9cm}}
\toprule
\textbf{Parámetro} & \textbf{Descripción} \\
\midrule
\texttt{query}      & Texto de la consulta a ejecutar. \\
\texttt{top\_k}     & Número máximo de resultados a devolver. \\
\texttt{source\_ids} & Lista opcional de identificadores de fuente para filtrar los
                       resultados a una o varias fuentes específicas. \\
\bottomrule
\end{tabular}
\end{center}

\section{Panel de configuración del pipeline}

\paragraph{Función del panel.}
Es la sección de mayor impacto operativo del sistema. Permite ajustar en tiempo de ejecución
todos los parámetros que controlan el comportamiento del pipeline RAG sin necesidad de
reiniciar el servidor ni modificar las variables de entorno. Los cambios se persisten en el
archivo \texttt{pipeline\_config.json} y sobreviven a reinicios del servidor, de modo que la
configuración ajustada se convierte en el nuevo estado base del sistema.

La existencia de este panel responde a la naturaleza exploratoria del ajuste de sistemas RAG.
Los valores óptimos de temperatura, número de fragmentos de contexto, pesos de fusión o
umbrales del detector de insuficiencia dependen del corpus específico y del perfil real de
consultas, y no son conocidos a priori. Poder modificarlos y observar el efecto sobre las
respuestas sin interrumpir el servicio reduce significativamente el ciclo de experimentación.

\paragraph{Datos presentados.}
El endpoint \texttt{GET /api/v1/config} devuelve la configuración activa completa. Al cargar
el panel se muestran todos los parámetros vigentes agrupados por categoría: pesos de fusión
híbrida, configuración del modelo LLM, estado y parámetros del reordenador, estado de HyDE,
y los quince parámetros del detector de insuficiencia.

\paragraph{Fusión híbrida.}

\begin{center}
\begin{tabular}{lll}
\toprule
\textbf{Parámetro} & \textbf{Default} & \textbf{Descripción} \\
\midrule
\texttt{bm25\_weight}   & 0.3 & Peso del recuperador léxico BM25 en la fusión RRF. \\
\texttt{vector\_weight} & 0.7 & Peso del recuperador vectorial en la fusión RRF. \\
\bottomrule
\end{tabular}
\end{center}

Los dos pesos deben sumar 1{,}0. Un peso BM25 más alto favorece coincidencias exactas de
términos técnicos; un peso vectorial más alto favorece similitud semántica entre idiomas
distintos. El sistema valida la restricción de suma antes de aplicar los cambios.

\paragraph{Modelo de lenguaje.}

\begin{center}
\begin{tabular}{lll}
\toprule
\textbf{Parámetro} & \textbf{Default} & \textbf{Descripción} \\
\midrule
\texttt{llm\_base\_url}  & \texttt{https://api.mistral.ai/v1} & URL del proveedor LLM. \\
\texttt{llm\_api\_key}   & ---                               & Clave API del proveedor. \\
\texttt{model}           & \texttt{mistral-small-latest}     & Modelo de generación. \\
\texttt{temperature}     & 0.1  & Temperatura de muestreo (0.0--2.0). \\
\texttt{max\_tokens}     & 1024 & Tokens máximos en la respuesta generada. \\
\texttt{context\_chunks} & 15   & Fragmentos incluidos en el prompt. \\
\bottomrule
\end{tabular}
\end{center}

\paragraph{Reordenador.}

\begin{center}
\begin{tabular}{lll}
\toprule
\textbf{Parámetro} & \textbf{Default} & \textbf{Descripción} \\
\midrule
\texttt{reranker\_enabled}      & true & Activa el cross-encoder de segunda etapa. \\
\texttt{reranker\_candidate\_k} & 30   & Candidatos pasados al reordenador. \\
\bottomrule
\end{tabular}
\end{center}

\paragraph{Expansión hipotética (HyDE).}

\begin{center}
\begin{tabular}{lll}
\toprule
\textbf{Parámetro} & \textbf{Default} & \textbf{Descripción} \\
\midrule
\texttt{hyde\_enabled} & false & Activa la expansión hipotética de la consulta antes de
                                 la búsqueda vectorial, a costa de una llamada LLM adicional. \\
\bottomrule
\end{tabular}
\end{center}

\paragraph{Detector de insuficiencia.}

\begin{center}
\begin{tabular}{lll}
\toprule
\textbf{Parámetro} & \textbf{Default} & \textbf{Descripción} \\
\midrule
\texttt{min\_results}                 & 5    & Mínimo de fragmentos para no activar búsqueda web. \\
\texttt{expected\_results}            & 10   & Resultados esperados (normaliza puntuación de cantidad). \\
\texttt{min\_top\_score}              & 0.35 & Puntuación mínima del fragmento más relevante. \\
\texttt{relevant\_overlap\_threshold} & 0.60 & Fracción mínima de términos de la consulta que debe
                                               contener un fragmento para ser considerado relevante. \\
\texttt{min\_relevant\_results}       & 2    & Mínimo de fragmentos relevantes requeridos. \\
\texttt{min\_coverage\_score}         & 0.20 & Cobertura léxica mínima de la consulta. \\
\texttt{min\_answerability\_score}    & 0.40 & Puntuación mínima de respondibilidad. \\
\texttt{min\_source\_diversity}       & 0.30 & Diversidad mínima de URLs únicas entre resultados. \\
\texttt{confidence\_threshold}        & 0.65 & Umbral de confianza local; por debajo activa la
                                               búsqueda web. \\
\texttt{coverage\_top\_n}             & 5    & Fragmentos del top usados para calcular cobertura. \\
\texttt{w\_top}                       & 0.10 & Peso de la puntuación máxima normalizada. \\
\texttt{w\_quantity}                  & 0.15 & Peso de la puntuación de cantidad. \\
\texttt{w\_coverage}                  & 0.35 & Peso de la cobertura léxica. \\
\texttt{w\_diversity}                 & 0.15 & Peso de la diversidad de fuentes. \\
\texttt{w\_answerability}             & 0.25 & Peso de la respondibilidad. \\
\bottomrule
\end{tabular}
\end{center}

Los cinco pesos (\texttt{w\_top}, \texttt{w\_quantity}, \texttt{w\_coverage},
\texttt{w\_diversity}, \texttt{w\_answerability}) deben sumar exactamente 1{,}0. La API
valida esta restricción antes de aplicar cualquier cambio y rechaza la configuración si no
se cumple.

\section{Panel de ingesta y fuentes}

\paragraph{Función del panel.}
Permite gestionar las fuentes de documentación del corpus: listar las fuentes configuradas,
ejecutar ingestas sobre URLs registradas y subir archivos directamente al sistema. Es la
sección desde la que el operador expande o actualiza el conocimiento disponible para el pipeline
RAG sin necesidad de acceder al servidor.

\paragraph{Datos presentados.}
El endpoint \texttt{GET /api/v1/sources} devuelve la lista de fuentes configuradas con su
identificador, URL raíz, estado y número de documentos indexados. El endpoint
\texttt{POST /api/v1/ingest} ejecuta la ingesta y devuelve el número de fragmentos procesados
y su distribución por fuente. El endpoint \texttt{POST /api/v1/upload} acepta archivos en
formato PDF, HTML, Markdown y texto plano; tras la carga informa el número de fragmentos
generados, los identificadores asignados y el estado de la operación.

\paragraph{Parámetros editables.}

\begin{center}
\begin{tabular}{lp{9cm}}
\toprule
\textbf{Parámetro} & \textbf{Descripción} \\
\midrule
\texttt{source\_id} & Identificador de la fuente a ingestar. \\
\texttt{force}      & Booleano que fuerza la re-ingesta aunque los documentos ya existan
                      en el índice. \\
\bottomrule
\end{tabular}
\end{center}

\section{Panel de métricas}

\paragraph{Función del panel.}
Muestra estadísticas operativas del sistema en tiempo real. Su utilidad es permitir al
operador monitorizar el estado del índice y el comportamiento del sistema sin acceder a los
logs del servidor. Es un panel de solo lectura; no expone parámetros editables.

\paragraph{Datos presentados.}
El endpoint \texttt{GET /api/v1/metrics} expone el número total de fragmentos en el índice
BM25 y en el índice vectorial, el número de documentos fuente indexados, el número de
fragmentos en el caché web, el tiempo de respuesta promedio de las últimas consultas y el
conteo de consultas que activaron búsqueda web externa.

\section{Panel de retroalimentación de consultas}

\paragraph{Función del panel.}
Implementa el módulo de \emph{query feedback}. Permite ejecutar búsquedas con expansión
pseudo-relevante de la consulta, registrar retroalimentación explícita sobre la relevancia
de los fragmentos recuperados, y realizar búsquedas que incorporan ese historial de
retroalimentación para reordenar los resultados. Es el panel desde el que el operador puede
mejorar la calidad de la recuperación de forma iterativa, sin modificar el índice ni el
pipeline.

\paragraph{Expansión de consulta.}
El endpoint \texttt{POST /api/v1/query-feedback/expand} aplica retroalimentación de
pseudo-relevancia: recupera los fragmentos más relevantes para la consulta, extrae los
términos más informativos y los añade a la consulta original. Devuelve la consulta expandida,
la lista de términos añadidos, el método de expansión utilizado y el número de fragmentos
empleados para la expansión.

\paragraph{Búsqueda con expansión y con retroalimentación.}
El endpoint \texttt{POST /api/v1/query-feedback/search} devuelve los campos de expansión más
una lista de resultados enriquecidos con el identificador, puntuación, URL, título, breadcrumb
y texto de cada fragmento. El endpoint \texttt{POST /api/v1/query-feedback/search-with-feedback}
incorpora adicionalmente el historial de retroalimentación explícita para ajustar las
puntuaciones. Cada resultado incluye la puntuación ajustada (\texttt{adjusted\_score}), el
boost aplicado (\texttt{feedback\_boost}), si se aplicó retroalimentación
(\texttt{feedback\_applied}), el valor de relevancia registrado (0--3), el tipo de coincidencia
con el historial (\texttt{exact} o \texttt{semantic}) y la similitud semántica con la consulta
original del registro.

\paragraph{Registro de retroalimentación.}
El endpoint \texttt{POST /api/v1/query-feedback/feedback} permite al usuario registrar su
valoración explícita sobre un fragmento recuperado. La interfaz muestra un selector de
relevancia junto a cada resultado de búsqueda. Las valoraciones posibles son: 0 (irrelevante),
1 (marginalmente relevante), 2 (relevante) y 3 (muy relevante). El campo \texttt{notes}
acepta un comentario libre opcional.

\paragraph{Parámetros editables.}

\begin{center}
\begin{tabular}{lll}
\toprule
\textbf{Parámetro} & \textbf{Tipo} & \textbf{Descripción} \\
\midrule
\texttt{query}                        & texto    & Consulta a ejecutar. \\
\texttt{top\_k}                       & entero   & Número de resultados (1--50). \\
\texttt{top\_k\_feedback}             & entero   & Fragmentos para la expansión (1--50). \\
\texttt{max\_expansion\_terms}        & entero   & Términos máximos de expansión (0--20). \\
\texttt{expansion\_enabled}           & booleano & Activa la expansión pseudo-relevante. \\
\texttt{feedback\_enabled}            & booleano & Activa el reordenamiento por retroalimentación. \\
\texttt{semantic\_feedback\_enabled}  & booleano & Activa la búsqueda semántica en el historial. \\
\texttt{semantic\_similarity\_threshold} & float & Umbral de similitud para coincidencias semánticas (0.0--1.0). \\
\texttt{chunk\_id}                    & texto    & Fragmento sobre el que se registra la valoración. \\
\texttt{relevance}                    & entero   & Valoración: 0 irrelevante, 1 marginal, 2 relevante,
                                                   3 muy relevante. \\
\texttt{notes}                        & texto    & Comentario libre opcional sobre la valoración. \\
\texttt{session\_id}                  & texto    & Identificador de sesión opcional. \\
\texttt{source\_ids}                  & lista    & Filtro opcional por fuentes específicas. \\
\bottomrule
\end{tabular}
\end{center}

\paragraph{Resumen de retroalimentación.}
El endpoint \texttt{GET /api/v1/query-feedback/feedback/summary} devuelve estadísticas
agregadas del historial: número total de valoraciones registradas, número de consultas únicas
con al menos una valoración, conteos de valoraciones positivas ($\geq 2$), negativas (= 0) y
marginales (= 1), y la media aritmética de todas las valoraciones.

\section{Panel de evaluación}

\paragraph{Función del panel.}
Permite gestionar el conjunto de evaluación del sistema: definir consultas de prueba con sus
juicios de relevancia (\emph{qrels}), ejecutar experimentos de recuperación y calcular métricas
de calidad estándar de recuperación de información. Es la sección desde la que el operador
puede medir el efecto de los cambios de configuración sobre la calidad objetiva del sistema.

\paragraph{Gestión de consultas y juicios.}
El endpoint \texttt{GET /api/v1/evaluation/queries} devuelve la lista de consultas de
evaluación registradas. El endpoint \texttt{POST /api/v1/evaluation/queries} permite añadir
nuevas consultas con su identificador, texto, descripción y juicios de relevancia asociados.
El endpoint \texttt{POST /api/v1/evaluation/queries/\{query\_id\}/rankings} ejecuta el
pipeline de recuperación para una consulta de evaluación y almacena el ranking resultante.
El endpoint \texttt{PUT /api/v1/evaluation/queries/\{query\_id\}/judgments/\{chunk\_id\}}
permite registrar o actualizar manualmente el juicio de relevancia de un fragmento específico.

\paragraph{Métricas calculadas.}
El endpoint \texttt{GET /api/v1/evaluation/summary} calcula y devuelve las métricas agregadas
del conjunto de evaluación completo: Precision@k, Recall@k, F1@k, MRR (\emph{Mean Reciprocal
Rank}), DCG@k y NDCG@k. El endpoint \texttt{GET /api/v1/evaluation/report} devuelve el reporte
detallado en formato JSON con métricas por consulta individual y promedios globales.

\paragraph{Parámetros editables.}

\begin{center}
\begin{tabular}{lp{9cm}}
\toprule
\textbf{Parámetro} & \textbf{Descripción} \\
\midrule
\texttt{query\_id} & Identificador de la consulta de evaluación a gestionar. \\
\texttt{chunk\_id} & Fragmento sobre el que se registra el juicio de relevancia. \\
\texttt{relevance} & Juicio de relevancia binario o graduado asignado al fragmento. \\
\bottomrule
\end{tabular}
\end{center}

\section{Resumen de endpoints por panel}

\begin{center}
\begin{tabular}{llp{6cm}}
\toprule
\textbf{Panel} & \textbf{Método} & \textbf{Endpoint} \\
\midrule
Chat RAG                  & POST & \texttt{/api/v1/rag/query} \\
Historial                 & GET  & \texttt{/api/v1/rag/history/\{session\_id\}} \\
Búsqueda híbrida          & POST & \texttt{/api/v1/search} \\
Búsqueda caché web        & POST & \texttt{/api/v1/search/web-cache} \\
Configuración (leer)      & GET  & \texttt{/api/v1/config} \\
Configuración (guardar)   & POST & \texttt{/api/v1/config} \\
Fuentes                   & GET  & \texttt{/api/v1/sources} \\
Ingesta                   & POST & \texttt{/api/v1/ingest} \\
Subida de archivo         & POST & \texttt{/api/v1/upload} \\
Métricas                  & GET  & \texttt{/api/v1/metrics} \\
Expansión de consulta     & POST & \texttt{/api/v1/query-feedback/expand} \\
Búsqueda con expansión    & POST & \texttt{/api/v1/query-feedback/search} \\
Búsqueda con feedback     & POST & \texttt{/api/v1/query-feedback/search-with-feedback} \\
Registrar valoración      & POST & \texttt{/api/v1/query-feedback/feedback} \\
Resumen de valoraciones   & GET  & \texttt{/api/v1/query-feedback/feedback/summary} \\
Valoraciones por consulta & GET  & \texttt{/api/v1/query-feedback/feedback} \\
Consultas de evaluación   & GET  & \texttt{/api/v1/evaluation/queries} \\
Nueva consulta            & POST & \texttt{/api/v1/evaluation/queries} \\
Ejecutar ranking          & POST & \texttt{/api/v1/evaluation/queries/\{id\}/rankings} \\
Juicio de relevancia      & PUT  & \texttt{/api/v1/evaluation/queries/\{id\}/judgments/\{cid\}} \\
Resumen de evaluación     & GET  & \texttt{/api/v1/evaluation/summary} \\
Reporte de evaluación     & GET  & \texttt{/api/v1/evaluation/report} \\
\bottomrule
\end{tabular}
\end{center}

\section{Justificación técnica de la solución}

La decisión de construir la interfaz como cliente puro de la API, sin lógica de negocio
propia, responde a un principio de mantenibilidad: el sistema tiene varios módulos con
parámetros ajustables en tiempo de ejecución, y cualquier duplicación de esa lógica en la
capa de presentación introduciría una fuente de inconsistencias. Al delegar toda la
computación al backend, la interfaz puede evolucionar independientemente del pipeline y
viceversa.

La organización en ocho paneles diferenciados refleja la separación de responsabilidades del
sistema. El panel de chat cubre el caso de uso primario; los paneles de búsqueda directa y
configuración cubren el uso operativo y de diagnóstico; los paneles de ingesta y métricas
cubren la administración del corpus; los paneles de retroalimentación y evaluación cubren el
ciclo de mejora continua de la calidad. Esta separación permite que distintos roles de usuario
(usuario final, operador, ingeniero de datos) interactúen únicamente con la sección que les
corresponde sin estar expuestos a la complejidad del resto.

La configuración en tiempo de ejecución expuesta en el panel de configuración es técnicamente
necesaria porque los parámetros del sistema RAG no tienen valores óptimos universales. El
umbral de confianza del detector de insuficiencia, los pesos de fusión híbrida o la temperatura
del modelo dependen del perfil real de consultas y del contenido del corpus. Sin la capacidad
de ajustarlos en caliente, cada experimento requeriría un reinicio del servidor y una
modificación de las variables de entorno, lo que hace imposible el ajuste interactivo en un
entorno de producción.

El panel de evaluación existe porque la calidad del sistema RAG no puede medirse únicamente
de forma cualitativa. Las métricas estándar de recuperación de información (NDCG, MRR,
Precision@k) proporcionan una señal objetiva y repetible que permite comparar configuraciones
distintas, detectar regresiones al cambiar modelos y justificar decisiones técnicas con
evidencia medible en lugar de impresiones subjetivas.

\section{Limitaciones del módulo}

La interfaz no tiene mecanismos de autenticación ni control de acceso. Cualquier usuario con
acceso al puerto del servidor puede modificar la configuración del pipeline, disparar
ingestas o registrar valoraciones. Esto es aceptable en un entorno de desarrollo interno pero
no en un despliegue con acceso no restringido.

El panel de historial solo puede recuperar sesiones por su identificador exacto. No existe
funcionalidad de búsqueda o listado de sesiones activas: el usuario debe conocer el
\texttt{session\_id} que usó en las consultas previas para poder acceder al historial. Sesiones
cuyo identificador se pierda son irrecuperables desde la interfaz, aunque los datos persistan
en Redis hasta que expire el TTL.

El panel de métricas expone estadísticas de nivel de sistema pero no métricas de calidad de
recuperación en tiempo real. Para conocer si el sistema está respondiendo bien a las consultas
recientes, el operador debe usar el panel de evaluación, que requiere construir un conjunto de
prueba y ejecutar los rankings manualmente.

Los cambios de configuración del detector de insuficiencia surten efecto inmediatamente sobre
las consultas siguientes pero no hay forma de comparar desde la interfaz el efecto de dos
configuraciones distintas sobre el mismo conjunto de consultas sin registrar manualmente los
resultados. La experimentación controlada requiere construir un conjunto de evaluación en el
panel de evaluación antes de realizar cambios.

\end{document}
